\begin{saresult}{A.121}{Proved under stated hypotheses}{A complete presentation for arbitrary PL manifolds in dimension at least nine}
Let $M$ be a compact connected PL manifold of dimension $n\geq9$,
possibly with boundary. All concordances fix the side boundary.
Choose a finite combinatorial triangulation $K$ of $M$, and let
$K^{(3)}$ be its $3$-skeleton. Then the countable abelian group
presentation $E(K^{(3)})$ of A.119, including every cycle,
boundary, and unit-image relation specified there, presents
\[
             \pi_0 P_{\rm PL}(M).
\]
No smoothing of $M$ is assumed. The same presentation gives the
components of the stipulated definable concordance group of any
compact definable manifold with this PL model, over any real closed
field and in any fixed o-minimal expansion. Thus this is an unstable
component computation, whereas A.120 computes the stable group in
every dimension. The presentation is countable; no finite invariant
factor decomposition or solution to its word problem is asserted.
For a disconnected compact manifold, use the finite product of the
presentations of its connected components.
\end{saresult}
\begin{proof}
We first work over $\mathbb R$. If $M$ has boundary, choose a PL collar
and push its outside face a positive distance into the collar. This
gives a PL embedding $e:M\longrightarrow\operatorname{int}M$
homotopic, as a map to $M$, to the identity. For example, on a collar
$\partial M\times[0,2]$ replace its second coordinate $t$ by
$\epsilon+(1-\epsilon/2)t$, with $0<\epsilon<1$, and leave the
complement of the collar fixed. The formulas agree at $t=2$;
replacing $\epsilon$ by $s\epsilon$ supplies the homotopy.
If $M$ is closed, take $e=\operatorname{id}$.
Put $L=e(|K^{(3)}|)$ and choose a sufficiently small closed regular
neighborhood $N$ of $L$ in $\operatorname{int}M$.

Here are the required properties of this neighborhood. They can all
be read from a finite subdivision, so no limiting neighborhood is
being substituted for $N$. Triangulate the pair $(M,L)$ with a mesh
small enough near $L$, and take its first derived subdivision, in
which the subdivision of $L$ is full.
If $\lambda$ is the sum of barycentric coordinates at its vertices,
the regular neighborhood may be taken to be $\lambda\geq1/2$ after
subdividing the cut cells. Shrinking inside the chosen open
neighborhood assures $N\subset\operatorname{int}M$. Its boundary
is $\lambda=1/2$, and it is a compact PL $n$-manifold. These are
the derived-neighborhood and simplicial-neighborhood constructions
of \cite[Chapter~3, Proposition~3.10 and Theorem~3.11, pp.~34--35,
and Corollary~3.18, p.~36]
{RourkeSanderson1972}.

Within each simplex write a point with $0<\lambda<1$ as
$\lambda u+(1-\lambda)v$, where $u$ uses its $L$-vertices and
$v$ its other vertices. Moving $\lambda$ to $1$ retracts $N$ onto
$L$. Moving it to $1/2$ retracts $N\setminus L$ onto $\partial N$.
The barycentric formulas agree on common faces, and are continuous
on their respective domains. In particular $N$ has the homotopy
type of the finite $3$-complex $K^{(3)}$ and is connected.

The inclusion $K^{(3)}\longrightarrow K$ is $3$-connected, since
all the omitted cells have dimension at least four. The retraction
and the homotopy from $e$ to the identity therefore show that
\[
                 \pi_i(M,N)=0\qquad (i\leq3).
                                                        \tag{A.121.1}
\]
Moreover, relative PL general position moves maps of intervals,
circles, and disks of dimension at most two off $L$, keeping any
part already in $N\setminus L$ fixed. The dimension inequality is
$2+3<n$. One may first push the portion of a map not fixed at
$\partial N$ into $\operatorname{int}N$ by a collar. The usual
relative general-position argument then applies there, and the
collar does not meet $L$; see \cite[Chapter~5, Theorems~5.3--5.4]
{RourkeSanderson1972}. For clarity in the injectivity argument,
represent a boundary loop by an embedded PL loop in $\partial N$,
which is possible by general position, and choose an embedded collar
on that loop into $N\setminus L$. Apply Theorem~5.4 to a filling
disk relative to this collar; its dimension bound makes the disk an
embedding. Then Theorem~5.3 moves its image off $L$ while fixing the
collar. Thus the relative data in the map theorem are indeed PL and
nondegenerate. The interval and circle arguments are the same in
lower dimension. Consequently
$N\setminus L\longrightarrow N$ induces isomorphisms on
$\pi_0$ and $\pi_1$. With the preceding boundary retraction this
proves
\[
       \pi_i(\partial N)\xrightarrow{\ \cong\ }\pi_i(N)
                      \qquad (i=0,1).                  \tag{A.121.2}
\]

The same regular-neighborhood construction gives $N$ a handle
decomposition with handles of indices at most three. Namely, thicken
the simplices of $L$ successively in increasing dimension. The
normal block of an $r$-simplex in the ambient PL manifold is a disk,
so the thickening step attaches $D^r\times D^{n-r}$ along
$\partial D^r\times D^{n-r}$. The product structures agree on
the previously thickened faces. This is the finite handle
construction from a triangulation; compare
\cite[Proposition~6.9 and its proof, pp.~82--83]
{RourkeSanderson1972}.
It also proves directly that $N$ has a compatible smooth structure.
Start with its $0$-handles and apply the classical PL--DIFF
handle-smoothing argument. The obstruction to extending over a
$k$-handle is the concordance class of the induced smoothing of
$S^{k-1}$; for $k=1,2,3$ this class is zero. The local product
structure and concordance-implies-isotopy theorems make the extension
agree with the already chosen smoothing near the attaching region.
This precise handle argument is given in
\cite[Essay~IV, Appendix~B, p.~211]{KirbySiebenmann1977}; it supplies
a smoothing compatible with the given PL structure, rather than
merely a smoothing of the underlying topological space. In
obstruction-theoretic language, the same conclusion follows by
lifting the PL tangent classifying map over a $3$-dimensional spine;
only the low, vanishing obstruction groups occur.

We now apply the PL localization theorem for concordances of
Burghelea--Lashof--Rothenberg. Its exact form is restated, with its
extension-by-identity map, in
\cite[Theorem~3.3, p.~17]{GalvinNonino2025}, citing
\cite[Theorem~3.1$'$]{BurgheleaLashofRothenberg1975}.
For $W\subset\operatorname{int}V$, it requires
$\pi_i(\partial W)\cong\pi_i(W)$ for $i=0,1$,
$\pi_i(V,W)=0$ for $i\leq r$, and $W$ to be $k$-connected,
where $0\leq k\leq r\leq\dim V-4$. It says that the map
\[
       \gamma:P_{\rm PL}(W)\longrightarrow P_{\rm PL}(V),
       \qquad \gamma(h)=h\text{ on }W\times I,
\]
extended by the identity elsewhere, has zero relative homotopy
groups in degrees at most
$\min(2r-3,r+k-2,r+2)$.
The relative pair is $(V,W)$: this agrees with the original
Theorem~3.1 on printed p.~28, which has been opened and visually
checked; the displayed $(W,V)$ in the later restatement reverses
the pair. The original pp.~28--29 establish the precursor theorem
and its handle proof. The later Theorem~3.1$'$ is used through the
explicit, opened restatement just cited; its original page has not
been retrieved.

Take $V=M$, $W=N$, $r=3$, and $k=0$. Equations
(A.121.1)--(A.121.2) check the hypotheses, and $3\leq n-4$.
The bound is
\[
                   \min(3,1,5)=1.
\]
Thus $\gamma$ induces a bijection on path components.
It preserves composition on the nose, because both extensions are
the identity off $N\times I$. Its component bijection is therefore
a group isomorphism
\[
          \pi_0P_{\rm PL}(N)\xrightarrow{\ \cong\ }
          \pi_0P_{\rm PL}(M).                          \tag{A.121.3}
\]
This step uses the actual localization map, rather than an unmarked
homotopy equivalence of the underlying spaces.

The manifold $N$ is smoothable, has dimension $n\geq9$, and has
finite homotopy model $K^{(3)}$. Apply A.119, including its boundary
version, to obtain
$\pi_0P_{\rm PL}(N)\congE(K^{(3)})$ as groups.
Together with (A.121.3) this proves the asserted full presentation.
It is unnecessary here to identify this presentation with
$E(K)$ or to assert a PL stabilization theorem for $M$.

Finally take a compact definable manifold over a real closed field
and a finite PL model $K$. Realize the same abstract triangulation
over the real algebraic numbers and hence over $\mathbb R$.
A.115 identifies the PL concordance components over these fields,
and the relative, component comparison A.113 identifies them with
the definable components in any fixed o-minimal expansion. Both
comparisons preserve composition. Thus the real group presentation
above gives exactly the stipulated definable component group as
well. No real model of the o-minimal expansion is used. A concordance
fixing the bottom cannot permute the finitely many connected
components of $M$, which proves the final product assertion.
\end{proof}
